% =============================================================================
% Review total de l'architecture du projet MyApp
% Compilation : pdflatex architecture-projet-myapp.tex
% =============================================================================

\documentclass[11pt,a4paper]{article}
\usepackage[utf8]{inputenc}
\usepackage[T1]{fontenc}
\usepackage[french]{babel}
\usepackage{geometry}
\usepackage{hyperref}
\usepackage{booktabs}
\usepackage{tabularx}
\usepackage{enumitem}
\usepackage{listings}
\usepackage{xcolor}
\usepackage{graphicx}
\usepackage{verbatim}
\usepackage{tikz}

\geometry{margin=2.5cm}
\hypersetup{colorlinks=true, linkcolor=blue, urlcolor=blue}
\setlength{\parindent}{0pt}
\setlength{\parskip}{0.6em}

% Style des listings
\lstset{
  basicstyle=\ttfamily\footnotesize,
  breaklines=true,
  frame=single,
  backgroundcolor=\color{gray!8},
  numbers=left,
  numberstyle=\tiny
}

\title{\textbf{Review total de l'architecture MyApp} \\ \large Projet, repos, cluster, CI/CD et déploiement}
\author{Équipe MyApp DevOps}
\date{\today}

\begin{document}

\maketitle
\tableofcontents
\newpage

%==============================================================================
\section{Introduction}
%==============================================================================

Ce document décrit une \textbf{review complète} de l'architecture du projet MyApp : structure des dépôts, services applicatifs, infrastructure Kubernetes, flux CI/CD et déploiement GitOps.

\textbf{Stack technique principale :}
\begin{itemize}[noitemsep]
  \item Cluster k3s (Kubernetes) sur GCP (2 nœuds)
  \item Namespace : \texttt{myapp}
  \item Ingress : Traefik (intégré k3s), IngressRoute pour le chatbot
  \item Service mesh / découverte : Consul
  \item GitOps : ArgoCD (branche \texttt{test-argocd})
  \item Accès externe : \texttt{https://example.ngrok-free.app}
\end{itemize}

%==============================================================================
\section{Structure des dépôts (code)}
%==============================================================================

Le projet est organisé en \textbf{monorepo virtuel} : plusieurs dépôts distincts, chacun avec son CI propre, contribuant au même produit MyApp.

\subsection{Liste des dépôts}

\begin{center}
\resizebox{\textwidth}{!}{%
\small
\setlength{\tabcolsep}{3pt}
\begin{tabular}{@{}llll@{}}
\toprule
\textbf{Dépôt} & \textbf{Stack} & \textbf{Branche CI} & \textbf{Image Docker} \\
\midrule
portable-platform-infra & Kustomize, K8s & test-argocd & (manifests) \\
Backend-api-gateway & Java / Spring Cloud Gateway & — & ghcr.io/.../backend-api-gateway \\
Backend-User-management-and-subscripption & Java / Spring Boot & test-ci-user-management & ghcr.io/.../backend-user-management... \\
Backend-payment-service & Java / Spring Boot & — & ghcr.io/.../backend-payment-service \\
Backend-crm-client & Java / Spring Boot & — & ghcr.io/.../backend-crm-client \\
Backend-predictions-intake-service & Java / Spring Boot & — & ghcr.io/.../backend-predictions-intake... \\
Backend-KPI-Dashboard-notifications & Java / Spring Boot & test-ci-kpi & ghcr.io/.../backend-kpi-dashboard... \\
Backend-chatbot & Python / FastAPI & test-ci-chatbot & ghcr.io/.../backend-chatbot \\
Front-end & Next.js & test-ci-frontend & ghcr.io/.../front-end \\
\bottomrule
\end{tabular}%
}
\end{center}

\subsection{Rôle du dépôt portable-platform-infra}

\begin{itemize}
  \item \textbf{Manifests Kubernetes} : base, infra, apps, cronjobs
  \item \textbf{ArgoCD} : Application pointant vers \texttt{deploy/k8s/}
  \item \textbf{Ansible} : provisionnement des VMs (k3s-server)
  \item \textbf{Docs} : TESTERS-GUIDE, CHECKLIST, DEPLOYMENT
\end{itemize}

%==============================================================================
\section{Architecture applicative}
%==============================================================================

\subsection{Vue d'ensemble des flux}

\begin{lstlisting}[language=,basicstyle=\ttfamily\small]
Ingress (Traefik) - host: example.ngrok-free.app
    /api, /login, /oauth2  -->  API Gateway (8888)
    /chatbot               -->  IngressRoute --> chatbot:8000
    /                      -->  Frontend (3000)

API Gateway -- Consul LB --> user-service, crmservice, payment-service,
                            prediction-intake-service, kpi-service
\end{lstlisting}

\subsection{Services backend}

\begin{center}
\begin{tabular}{@{}lllll@{}}
\toprule
\textbf{Service} & \textbf{Port} & \textbf{Nom Consul} & \textbf{Auth} & \textbf{Réplicas} \\
\midrule
api-gateway       & 8888 & api-gateway        & —      & 2 \\
user-management   & 8081 & user-service       & JWT/OAuth & 2 \\
crm-client        & 8083 & crmservice         & JWT    & 2 \\
payment-service   & 8082/8083 & payment-service & JWT   & 2 \\
predictions-intake& 8082 & prediction-intake-service & JWT & 2 \\
kpi-dashboard     & 8084 & kpi-service        & JWT    & 1 \\
chatbot           & 8000 & (hors gateway)     & Non    & 1 \\
frontend          & 3000 & —                  & —      & 2 \\
\bottomrule
\end{tabular}
\end{center}

\subsection{Routing API Gateway (Spring Cloud Gateway)}

\begin{center}
\begin{tabular}{@{}p{6cm}p{5cm}@{}}
\toprule
\textbf{Préfixe} & \textbf{Service cible} \\
\midrule
\texttt{/api/auth/**}, \texttt{/auth/**}, \texttt{/oauth2/**} & user-service \\
\texttt{/api/v1/predictions/**} & prediction-intake-service \\
\texttt{/api/v1/api-keys/**} & user-service \\
\texttt{/api/kpi/**} & kpi-service \\
\texttt{/internal/crm/**} & crmservice \\
\texttt{/payment-service/internal/payments/**} & payment-service \\
\texttt{/payment-service/api/v1/webhooks/**} & payment-service (Stripe) \\
\texttt{/chatbot/**} & (IngressRoute Traefik direct) \\
\bottomrule
\end{tabular}
\end{center}

\subsection{Chatbot (FastAPI)}

\begin{itemize}
  \item \textbf{Stack} : Python 3.11, FastAPI, OpenRouter (LLM), Knowledge Base locale
  \item \textbf{Accès} : IngressRoute Traefik avec \texttt{stripPrefix} sur \texttt{/chatbot}
  \item \textbf{Endpoints} : \texttt{/health}, \texttt{/chat}
  \item \textbf{Secret} : \texttt{chatbot-credentials} (LLM\_API\_KEY) créé manuellement
\end{itemize}

%==============================================================================
\section{Architecture infrastructure Kubernetes}
%==============================================================================

\subsection{Structure Kustomize}

\begin{verbatim}
deploy/k8s/
|-- base/           # namespace myapp
|-- infra/          # PostgreSQL, Redis, Consul, RabbitMQ
|   |-- postgres/   # PVC 5Gi (local-path), secret
|   |-- redis/      # 2 replicas
|   |-- consul/     # 1 replica (bootstrap-expect=1)
|   |-- rabbitmq/   # 2 replicas
|-- apps/           # api-gateway, frontend, user-management, ...
|   |-- ingress-ip.yaml   # ngrok + IP
|   |-- chatbot/          # IngressRoute + Middleware stripPrefix
|   +-- ...
|-- cronjobs/       # clean-disk-backend, clean-disk-frontend, clean-evicted-pods
|-- scripts/        # clean-node-disk.sh
\end{verbatim}
\end{verbatim}

\subsection{Nœuds}

\begin{center}
\begin{tabular}{@{}lllll@{}}
\toprule
\textbf{Nœud} & \textbf{Rôle} & \textbf{IP interne} & \textbf{CPU} & \textbf{RAM} \\
\midrule
backend-vm  & control-plane, master & 10.0.0.11 & $\sim$2 cores & $\sim$4 Go \\
frontend-vm & worker                 & 10.0.0.12 & $\sim$2 cores & $\sim$4 Go \\
\bottomrule
\end{tabular}
\end{center}

\subsection{Stockage}

\begin{itemize}
  \item \textbf{PostgreSQL} : PVC 5Gi (local-path), RWO
  \item Pas de PVC pour Redis, RabbitMQ (ephemeral)
\end{itemize}

\subsection{Secrets (Kubernetes)}

\begin{center}
\begin{tabular}{@{}ll@{}}
\toprule
\textbf{Secret} & \textbf{Usage} \\
\midrule
ghcr-secret & Pull images depuis GHCR \\
postgres-credentials & PostgreSQL \\
chatbot-credentials & LLM\_API\_KEY (OpenRouter) \\
stripe-credentials & Payment \\
google-oauth-credentials & OAuth Google \\
mail-credentials & Email \\
user-management secret & JWT, Google, etc. \\
\bottomrule
\end{tabular}
\end{center}

\subsection{CronJobs}

\begin{center}
\begin{tabular}{@{}lll@{}}
\toprule
\textbf{CronJob} & \textbf{Schedule} & \textbf{Rôle} \\
\midrule
clean-disk-frontend  & 15,45 * * * * & Libérer disque frontend-vm \\
clean-disk-backend  & */30 * * * *  & Libérer disque backend-vm \\
clean-evicted-pods  & */30 * * * *  & Supprimer pods Evicted \\
\bottomrule
\end{tabular}
\end{center}

%==============================================================================
\section{CI/CD et déploiement GitOps}
%==============================================================================

\subsection{Flux global}

\begin{enumerate}
  \item \textbf{CI par service} : push sur branche test-ci-* $\rightarrow$ build Docker $\rightarrow$ push GHCR
  \item \textbf{Update manifests} : le workflow clone portable-platform-infra, met à jour le tag image dans le deployment correspondant, push sur \texttt{test-argocd}
  \item \textbf{ArgoCD} : observe \texttt{test-argocd}, \texttt{syncPolicy.automated} + \texttt{selfHeal}
  \item \textbf{Résultat} : déploiement automatique sans intervention manuelle
\end{enumerate}

\subsection{Prérequis CI : GH\_PAT}

Chaque dépôt qui met à jour les manifests doit avoir le secret \texttt{GH\_PAT} (Personal Access Token avec scope \texttt{repo}) pour cloner et pousser vers portable-platform-infra.

\subsection{Exemple CI (Backend-chatbot)}

\begin{lstlisting}
# Branches : test-ci-chatbot
# 1. Checkout
# 2. Python 3.11, pip install, pytest
# 3. Docker build (Dockerfile.optimized)
# 4. Push ghcr.io/myapp/backend-chatbot:${SHA:0:7}
# 5. Update deploy/k8s/apps/chatbot/deployment.yaml
# 6. Commit + push vers test-argocd
\end{lstlisting}

\subsection{ArgoCD Application}

\begin{lstlisting}
source:
  repoURL: https://github.com/MyApp/portable-platform-infra.git
  targetRevision: test-argocd
  path: deploy/k8s
destination:
  server: https://kubernetes.default.svc
  namespace: myapp
syncPolicy:
  automated:
    prune: true
    selfHeal: true
\end{lstlisting}

%==============================================================================
\section{Ingress et accès réseau}
%==============================================================================

\subsection{Ingress principaux}

\begin{center}
\begin{tabular}{@{}lll@{}}
\toprule
\textbf{Type} & \textbf{Host / Path} & \textbf{Backend} \\
\midrule
Ingress & \texttt{api.localhost} & api-gateway:8888 \\
Ingress & \texttt{app.localhost} & frontend:3000 \\
Ingress & \texttt{chatbot.localhost} & chatbot:8000 \\
IngressRoute & \texttt{chatbot-ngrok} & chatbot + stripPrefix \texttt{/chatbot} \\
Ingress (ngrok) & \texttt{example.ngrok-free.app/api} & api-gateway \\
Ingress (ngrok) & \texttt{.../chatbot} & (via IngressRoute) \\
Ingress (ngrok) & \texttt{.../} & frontend \\
\bottomrule
\end{tabular}
\end{center}

\subsection{Priorité Ingress (ngrok)}

\begin{itemize}
  \item \texttt{ingress-ip-api} : priorité 200 (\texttt{/api}, \texttt{/login}, \texttt{/oauth2})
  \item \texttt{ingress-ip-frontend} : priorité 100 (\texttt{/} catch-all)
\end{itemize}

\subsection{Header ngrok}

Toutes les requêtes API manuelles doivent inclure :
\begin{lstlisting}
ngrok-skip-browser-warning: 1
\end{lstlisting}

%==============================================================================
\section{Images Docker et registre}
%==============================================================================

\begin{center}
\begin{tabular}{@{}ll@{}}
\toprule
\textbf{Service} & \textbf{Image (ghcr.io/myapp/...)} \\
\midrule
api-gateway & backend-api-gateway \\
frontend & front-end \\
user-management & backend-user-management-and-subscripption \\
payment-service & backend-payment-service \\
crm-client & backend-crm-client \\
predictions-intake & backend-predictions-intake-service \\
kpi-dashboard & backend-kpi-dashboard-notifications \\
chatbot & backend-chatbot \\
\bottomrule
\end{tabular}
\end{center}

Tags : généralement le SHA court du commit (ex. \texttt{555848d}).

%==============================================================================
\section{Points de review et recommandations}
%==============================================================================

\subsection{Points positifs}

\begin{itemize}
  \item Architecture claire : API Gateway central, services enregistrés dans Consul
  \item GitOps avec ArgoCD : déploiements reproductibles
  \item CI/CD par service avec mise à jour automatique des manifests
  \item HA sur les services critiques (2 replicas pour gateway, frontend, redis, rabbitmq)
  \item CronJobs pour maintenir la santé du cluster (disque, pods evicted)
  \item Documentation : TESTERS-GUIDE, INDEX, CHECKLIST
  \item Chatbot isolé du gateway (IngressRoute directe), stripPrefix correct
\end{itemize}

\subsection{Points d'attention}

\begin{enumerate}
  \item \textbf{Consul} : 1 seule instance $\rightarrow$ pas de HA. En prod, envisager 3 nodes.
  \item \textbf{PostgreSQL} : 1 replica, PVC local-path $\rightarrow$ risque de perte de données si le nœud tombe.
  \item \textbf{Secret chatbot} : \texttt{chatbot-credentials} doit être créé manuellement (LLM\_API\_KEY).
  \item \textbf{Mémoire} : frontend-vm souvent chargée ($\sim$83\% RAM). Surveiller.
  \item \textbf{HTTPS} : accès via ngrok (HTTPS côté ngrok), pas de certificats internes.
  \item \textbf{Service user-service} : service K8s distinct de user-management $\rightarrow$ vérifier la cohérence.
\end{enumerate}

\subsection{Recommandations}

\begin{itemize}
  \item Continuer à utiliser des tags fixes (SHA) plutôt que \texttt{latest}
  \item En production : renforcer secrets (Sealed Secrets, Vault)
  \item Configurer \texttt{FRONTEND\_URL} en variable d'environnement (OAuth callback)
  \item JWT\_SECRET partagé entre api-gateway et user-management
\end{itemize}

%==============================================================================
\section{Références}
%==============================================================================

\begin{itemize}
  \item \texttt{deploy/docs/INDEX.md} — index des documentations
  \item \texttt{deploy/TESTERS-GUIDE.md} — guide testeurs, APIs, pentest
  \item \texttt{deploy/k8s/README.md} — structure K8s
  \item \texttt{deploy/k8s/CHECKLIST.md} — checklist configuration
  \item \texttt{deploy/argocd/ARGOCD-AUTODEPLOY.md} — setup GH\_PAT
  \item \texttt{deploy/docs/architecture-cluster-myapp.tex} — doc cluster détaillée
\end{itemize}

\vspace{1em}
\hrule
\vspace{0.5em}
\textit{Review architecture MyApp — Document généré pour l'équipe DevOps.}

\end{document}
